% from ega4-15-fr.tex
\oldpage[IV-3]{239}
en outre $Y$ noethérien~; comme $f$ est universellement ouvert (11.1.1), il suffit alors d'appliquer (15.6.4).
\end{proof}

\begin{proposition}[15.6.6]
\EGAReaderTarget{ega4.c04.prop.15-6-6}{IV.15.6.6-fr}
Soit $f:X\to Y$ un morphisme tel que $Y$ soit localement noethérien et $f$ localement de type fini, ou que $f$ soit de présentation finie. Soit $g$ une $Y$-section de $X$~; pour tout $y\in Y$, soit $X_y^0$ la composante connexe de $f^{-1}(y)$ contenant $g(y)$~; enfin, soit $X^0$ la réunion des $X_y^0$ lorsque $y$ parcourt $Y$.

Supposons que, pour tout $y\in Y$, $X_y^0$ soit géométriquement irréductible. Alors :
\begin{enumerate}[label=\emph{(\roman*)}]
  \item La fonction $y\mapsto\dim(X_y^0)$ est semi-continue supérieurement dans $Y$.
  \item Soient $x_0\in X^0$, $y_0=f(x_0)$. Si la fonction $y\mapsto\dim(X_y^0)$ est constante dans un voisinage de $y_0$, il existe un voisinage ouvert $U$ de $y_0$ tel que $f$ soit universellement ouvert aux points de $X^0\cap f^{-1}(U)$. Si de plus $Y$ est réduit et $f^{-1}(y_0)$ géométriquement réduit sur $\mathbf{k}(y_0)$ au point $x_0$, $f$ est plat au point $x_0$.
  \item Inversement, supposons que $f$ soit universellement ouvert au point générique de $X_{y_0}^0$ et en outre que l'une des conditions suivantes soit vérifiée :
  \begin{enumerate}
    \item[$\alpha$)] La fibre $f^{-1}(y_0)$ est géométriquement réduite sur $\mathbf{k}(y_0)$ au point $g(y_0)$ et l'anneau $\mathcal O_{Y,y_0}$ est noethérien.
    \item[$\beta$)] Pour tout point générique $y'$ d'une composante irréductible de $Y$ contenant $y_0$, on a $\dim(f^{-1}(y'))=\dim(X_{y'}^0)$.
  \end{enumerate}
  Alors la fonction $y\mapsto\dim(X_y^0)$ est constante dans un voisinage de $y_0$.
  \item L'ensemble $W$ des points $x\in X^0$ tels que la fonction $y\mapsto\dim(X_y^0)$ soit constante dans un voisinage de $f(x)$ et que $X_{f(x)}^0$ soit géométriquement réduit sur $\mathbf{k}(f(x))$, est ouvert dans $X$.
\end{enumerate}
\end{proposition}

\begin{proof}
Les questions étant locales sur $Y$, on peut supposer que $Y=\operatorname{Spec}(A)$ est affine. Lorsque $f$ est de présentation finie, il existe donc un sous-anneau noethérien $A_1$ de $A$ et un morphisme de type fini $f_1:X_1\to Y_1=\operatorname{Spec}(A_1)$ tels que $X=X_1\times_{Y_1}Y$ et $f=(f_1)_{(Y)}$ (8.9.1)~; en outre, on peut supposer qu'il existe une $Y_1$-section $g_1$ de $X_1$ telle que $g=(g_1)_{(Y)}$ (8.8.2). Pour tout $y\in Y$, soit $y_1$ sa projection dans $Y_1$~; il résulte de (4.5.13) que $(X_1)_{y_1}^0$ est géométriquement connexe, donc, si $p:f^{-1}(y)\to f_1^{-1}(y_1)$ est la projection canonique, on a $X_y^0=p^{-1}((X_1)_{y_1}^0)$ (4.4.1). Notons d'autre part que l'on peut, en vertu de (9.7.11) et (9.7.12), appliquer (9.3.3) à la propriété d'être géométriquement irréductible~; on peut donc supposer $A_1$ choisi de telle sorte que $(X_1)_{y_1}^0$ soit géométriquement irréductible pour tout $y_1\in Y_1$. On a $\dim(X_y^0)=\dim((X_1)_{y_1}^0)$ (4.1.4), et en appliquant de nouveau (9.3.3) à la propriété d'avoir une dimension donnée (et utilisant (9.5.5) et (9.7.12)), on peut supposer (en restreignant au besoin $Y$ à un voisinage de $y_0$) que si $\dim(X_y^0)$ est constante dans $Y$, alors $\dim((X_1)_{y_1}^0)$ est constante dans $Y_1$. Si $Y$ est réduit, il en est de même de $Y_1$ puisque $A_1\subset A$~; d'autre part, si $f^{-1}(y_0)$ est géométriquement réduit au point $x_0$, $f_1^{-1}(y_{01})$ l'est au point $x_{01}$ ($x_{01}$ et $y_{01}$ étant les projections respectives dans $X_1$ et $Y_1$ de $x_0$ et $y_0$) (4.6.10).

Ces remarques montrent que pour démontrer (i) et (ii), on peut se borner au cas où $Y$ est noethérien et $f$ localement de type fini. Mais dans ce cas, l'assertion (i) résulte de (15.6.2) et la première assertion de (ii) de (15.6.3). D'autre part, si $Y$ est réduit et $f^{-1}(y_0)$ géométriquement réduit sur $\mathbf{k}(y_0)$ au point $x_0$ ($Y$ étant toujours supposé

